%% kapa: exam=1 round=34 year=2023 period=1 form=A subject=real_estate subject_ko=부동산학원론 seq=37 number=117
%% source: Q-Net 감정평가사 기출문제 (cst00309) artlSeq=5211789 file=raw/1008/5211789/2023년도 제34회 감정평가사 1차시험 1교시.pdf
%% conversion: kordoc 4.24.51 -> tools/gen1.py (auto); review=proofread-sonnet; flags=none
\begin{question}{117}
감정평가방식 중 원가방식에 관련된 설명으로 옳은 것은?
\begin{choices}
  \item 원가방식은 대체의 원칙, 수요와 공급의 원칙, 균형의 원칙, 외부의 원칙, 예측의 원칙과 밀접한 관련이 있다.
  \item 재조달원가란 대상물건을 기준시점에 재생산 또는 재취득하는데 필요한 적정원가의 총액으로서 원칙적으로 그 대상물건 값의 상한선을 나타낸다.
  \item 대치원가(replacement cost)란 건축자재, 설비공법 등에 있어 신축시점의 표준적인 것을 사용한 적정원가로서 이미 기능적 감가는 반영되어 있다.
  \item 재조달원가를 구하는 방법은 직접법으로 총가격적산법(총량조사법), 변동률적용법(비용지수법) 등이 있고, 간접법으로 부분별단가적용법, 단위비교법 등이 있다.
  \item 감가수정에 있어서 감가요인은 물리적요인, 기능적요인, 경제적요인이 있으며, 감가상각에 있어서 감가요인은 물리적요인, 경제적요인이 있다.
\end{choices}
\end{question}
\answer{2}
